% ============================================================
% COMPACT POST-PROCESSING VISUALIZATION
% Cleaner, more compact version for document integration
% ============================================================

\begin{figure}[htbp]
\centering
\begin{tikzpicture}[
    scale=0.95,
    node distance=1.5cm,
    every label/.append style={font=\sffamily\tiny}
]

% Colors
\definecolor{algo1}{RGB}{230, 126, 34}
\definecolor{algo2}{RGB}{41, 128, 185}
\definecolor{algo3}{RGB}{39, 174, 96}
\definecolor{algo4}{RGB}{155, 89, 182}
\definecolor{nodecolor}{RGB}{52, 73, 94}
\definecolor{bccolor}{RGB}{155, 89, 182}
\definecolor{removedcolor}{RGB}{231, 76, 60}

% Styles
\tikzset{
    node_free/.style={circle, minimum size=6pt, fill=nodecolor, draw=nodecolor, line width=1pt},
    node_bc/.style={circle, minimum size=7pt, fill=bccolor, draw=bccolor, line width=1.2pt},
    edge/.style={line width=1.2pt, draw=black!40},
    edge_var/.style={line width=1.2pt},
    algo_box/.style={
        rectangle,
        rounded corners=2pt,
        minimum width=2.5cm,
        minimum height=0.8cm,
        text centered,
        text width=2.3cm,
        font=\sffamily\small\bfseries,
        draw,
        line width=1.5pt
    },
    algo1box/.style={algo_box, fill=algo1!20, draw=algo1!70, text=algo1!80},
    algo2box/.style={algo_box, fill=algo2!20, draw=algo2!70, text=algo2!80},
    algo3box/.style={algo_box, fill=algo3!20, draw=algo3!70, text=algo3!80},
    algo4box/.style={algo_box, fill=algo4!20, draw=algo4!70, text=algo4!80},
    step_label/.style={font=\sffamily\scriptsize, text centered, text width=2.2cm},
}

% ============================================================
% STEP 1: EDGE COLLAPSE
% ============================================================
\node[algo1box] (a1) at (0, 2) {Edge\\Collapse};
\node[step_label, below=0.6cm of a1] (a1_param) {$d_{\text{collapse}}$};

\node[node_bc] (n1_a) at (-0.4, 0.8) {};
\node[node_free] (n1_b) at (0, 1) {};
\node[node_free] (n1_c) at (0.4, 0.8) {};
\node[node_bc] (n1_d) at (0.8, 1) {};

\draw[edge] (n1_a) -- (n1_b);
\draw[edge, line width=0.8pt] (n1_b) -- (n1_c); % short edge to be collapsed
\draw[edge] (n1_c) -- (n1_d);

\node[font=\sffamily\tiny, below=1.8cm of a1, text centered, text width=2.3cm] {
    \small \textit{Merge short}\\
    \small \textit{edges}
};

% ============================================================
% STEP 2: BRANCH PRUNING
% ============================================================
\node[algo2box] (a2) at (3.5, 2) {Branch\\Pruning};
\node[step_label, below=0.6cm of a2] (a2_param) {$l_{\text{prune}}$};

\node[node_bc] (n2_a) at (3.1, 0.8) {};
\node[node_free] (n2_b) at (3.5, 1) {};
\node[node_free] (n2_c) at (3.9, 0.8) {};
\node[node_bc] (n2_d) at (4.3, 1) {};
\node[node_free] (n2_br) at (3.5, 1.5) {}; % branch to prune

\draw[edge] (n2_a) -- (n2_b);
\draw[edge] (n2_b) -- (n2_c);
\draw[edge] (n2_c) -- (n2_d);
\draw[edge, densely dotted, line width=0.8pt] (n2_b) -- (n2_br); % branch removed

\node[font=\sffamily\tiny, below=1.8cm of a2, text centered, text width=2.3cm] {
    \small \textit{Remove}\\
    \small \textit{weak branches}
};

% ============================================================
% STEP 3: RDP SIMPLIFICATION
% ============================================================
\node[algo3box] (a3) at (7, 2) {RDP\\Simplification};
\node[step_label, below=0.6cm of a3] (a3_param) {$\varepsilon_{\text{rdp}}$};

\node[node_bc] (n3_a) at (6.6, 0.8) {};
\node[node_free] (n3_b) at (7, 1) {};
\node[node_free] (n3_c) at (7.4, 0.8) {};
\node[node_bc] (n3_d) at (7.8, 1) {};
\node[node_free] (n3_m1) at (6.8, 0.95) {}; % intermediate point
\node[node_free] (n3_m2) at (7.2, 0.85) {}; % intermediate point

\draw[edge, dotted, line width=0.8pt] (n3_a) -- (n3_m1) -- (n3_m2) -- (n3_b); % many intermediate points
\draw[edge] (n3_b) -- (n3_c);
\draw[edge] (n3_c) -- (n3_d);

\node[font=\sffamily\tiny, below=1.8cm of a3, text centered, text width=2.3cm] {
    \small \textit{Reduce polyline}\\
    \small \textit{points}
};

% ============================================================
% STEP 4: RADIUS ESTIMATION (SINGLE BOX)
% ============================================================
\node[algo4box] (a4) at (10.5, 2) {Radius\\Estimation};

\node[step_label, below=0.6cm of a4, text width=2.8cm] (a4_param) {
    EDT: $r_v = \text{EDT}(v) \cdot h$ \\[2pt]
    Uniform: $r = \sqrt{\frac{V_t}{\pi\sum L_i}}$
};

\node[node_bc] (n4_a) at (10.1, 0.8) {};
\node[node_free] (n4_b) at (10.5, 1) {};
\node[node_free] (n4_c) at (10.9, 0.8) {};
\node[node_bc] (n4_d) at (11.3, 1) {};

% Show variable thickness (EDT)
\draw[edge_var, line width=2.2pt, draw=algo4] (n4_a) -- (n4_b);
\draw[edge_var, line width=1.8pt, draw=algo4] (n4_b) -- (n4_c);
\draw[edge_var, line width=2.4pt, draw=algo4] (n4_c) -- (n4_d);

\node[font=\sffamily\tiny, below=1.8cm of a4, text centered, text width=2.3cm] {
    \small \textit{Assign radii}\\
    \small \textit{(variable or constant)}
};

% ============================================================
% CONNECTING ARROWS
% ============================================================
\draw[->, line width=2pt, color=black!50] (a1) -- (a2);
\draw[->, line width=2pt, color=black!50] (a2) -- (a3);
\draw[->, line width=2pt, color=black!50] (a3) -- (a4);

% ============================================================
% LEGEND
% ============================================================
\node[rectangle, draw=black!25, line width=0.8pt, fill=white!90, inner sep=6pt,
      text width=10.5cm, font=\sffamily\scriptsize, below=2.2cm of a3] {
    \textbf{Legend:} \,
    \tikz[baseline=(n.base)]{\node[node_bc, inner sep=1.2pt] (n) {};}  BC node \quad
    \tikz[baseline=(n.base)]{\node[node_free, inner sep=1.2pt] (n) {};}  Free node \quad
    Edge thickness represents radius (Algorithm~4 only)
};

\end{tikzpicture}

\caption{\textbf{Four sequential post-processing algorithms.}
(1)~Edge Collapse: merge nodes connected by edges shorter than $d_{\mathrm{collapse}}$.
(2)~Branch Pruning: remove degree-1 branches with length $< l_{\mathrm{prune}}$, protecting BC nodes.
(3)~RDP Simplification: reduce intermediate polyline points with tolerance $\varepsilon_{\mathrm{rdp}}$.
(4)~Radius Estimation: assign cross-section radii via EDT (variable, geometry-based) or Uniform mode (constant, volume-conserving).}
\label{fig:postprocessing}
\end{figure}
